求める軌跡を$R$とし、中点$M$の座標を$M(x,y)$とする。
また、交点$A,B$の$x$座標をそれぞれ$\alpha,\beta$とする。

直線の式を双曲線の式へ代入すると、交点の$x$座標は
\[
8X^2+6kX+k^2+2=0
\]
を満たす。
この2次方程式の判別式を$D$とすると、
\[
D=(6k)^2-4\cdot8(k^2+2)
  =4(k^2-16)
\]
である。また、解と係数の関係から、
\[
\alpha+\beta=-\frac{3k}{4}
\]
である。

\[
\begin{aligned}
M(x,y)\in R
&\Longleftrightarrow
\begin{gathered}
\text{「直線 }y=3x+k\text{ が双曲線 }x^2-y^2=2\text{ と異なる2点 }A,B\text{ で交わり、}\\
\text{その中点が }M(x,y)\text{ となる実数 }k\text{ が存在する」}
\end{gathered}\\
&\Longleftrightarrow
\begin{gathered}
\text{「}\;
\left\{
\begin{aligned}
D&>0\\
x&=\frac{\alpha+\beta}{2}\\
y&=\frac{3(\alpha+\beta)+2k}{2}
\end{aligned}
\right.\\
\text{を満たす実数 }k\text{ が存在する」}
\end{gathered}\\
&\Longleftrightarrow
\begin{gathered}
\text{「}\;
\left\{
\begin{aligned}
|k|&>4\\
x&=-\frac{3k}{8}\\
y&=-\frac{k}{8}
\end{aligned}
\right.\\
\text{を満たす実数 }k\text{ が存在する」}
\end{gathered}\\
&\Longleftrightarrow
\left\{
\begin{aligned}
x&=3y\\
|x|&>\frac32
\end{aligned}
\right.
\end{aligned}
\]

したがって、求める軌跡は、直線$x=3y$のうち
$x<-\frac32$または$x>\frac32$を満たす部分である。
